% SPDX-License-Identifier: Apache-2.0
%
% What has been built: two pages on the language, the parts, the core
% and the GPU. Built by //docs:showcase. Every listing is included from
% the tree and the picture is the one the GPU drew.
\documentclass[journal,9pt]{IEEEtran}

\newcommand{\listingsize}{\fontsize{6}{7}\selectfont}
% SPDX-License-Identifier: Apache-2.0
%
% The house preamble, shared by every document in this directory. Loaded
% after \documentclass, before \title. Keeping it in one file is what
% stops two articles drifting apart in font, listing style or figure
% style.
% Computer Modern. IEEEtran selects Times (ptm) by default, so the face has
% to be chosen explicitly.
%
% Latin Modern is the Computer Modern variety used here, and the choice is
% forced rather than aesthetic. Under T1 encoding, plain `cmr` has no Type 1
% outlines unless cm-super is installed, and it is not in the pinned
% distribution, so pdflatex falls back to EC bitmaps and every glyph in the
% PDF becomes a Type 3 font. That was measured: the first build produced a
% document with 10 Type 3 fonts and no Type 1. Latin Modern is Computer
% Modern redrawn with a full T1 repertoire and real .pfb outlines, so the
% same design comes out scalable.
\usepackage[T1]{fontenc}
\usepackage[utf8]{inputenc}
\usepackage{lmodern}
% microtype is not in the pinned TeX distribution. emergencystretch alone
% keeps the two-column measure from producing overfull lines.
\setlength{\emergencystretch}{3em}

\usepackage{amsmath}
\usepackage{amssymb}
\usepackage{amsthm}
\usepackage{booktabs}
\usepackage{array}
% enumitem is not in the pinned TeX distribution; the list spacing below
% does what \setlist would have done.
\usepackage{float}
\usepackage{xcolor}
\usepackage{listings}
\usepackage{tikz}
\usetikzlibrary{arrows.meta,positioning,fit,backgrounds,calc}
\usepackage[hidelinks,bookmarks=true]{hyperref}

% Numbered, definition-style box for a rule the reader is meant to apply.
\theoremstyle{definition}
\newtheorem{rulebox}{Rule}
\newtheorem{example}{Example}

% Unnumbered asides.
\theoremstyle{remark}
\newtheorem*{tip}{Tip}
\newtheorem*{pitfall}{Pitfall}

% Tighter lists, in place of enumitem's \setlist.
\setlength{\topsep}{2pt}
\setlength{\itemsep}{1pt}
\setlength{\parsep}{0pt}

\newcommand{\code}[1]{\texttt{#1}}
\newcommand{\term}[1]{\emph{#1}}

% Syntax highlighting. Four roles, distinguishable in colour and still
% distinguishable in greyscale, because an IEEE article gets printed: the
% keyword is the only bold role, the comment is the only italic one, and the
% two remaining roles differ in lightness as well as in hue.
\definecolor{kwblue}{RGB}{20,60,150}
\definecolor{cmtgreen}{RGB}{90,120,90}
\definecolor{strred}{RGB}{150,50,40}
\definecolor{numgray}{gray}{0.45}
\definecolor{framegray}{gray}{0.70}
\definecolor{bgshade}{gray}{0.97}
% A darker fill for figure cells. The listing background has to stay very
% light to keep code readable; a figure cell has to be clearly shaded or the
% distinction it draws is invisible in print.
\definecolor{cellshade}{gray}{0.90}

% The listing size is the document's choice, because it depends on the
% column: a two-column article sets `\listingsize` to 6pt and keeps its
% listings within 70 columns, a one-column document keeps the body size
% and includes source files at 80. Lines are never broken; a line that does not fit
% overflows the frame, which is visible, rather than wrapping, which is
% not.
\providecommand{\listingsize}{\scriptsize}

% `'` is deliberately not a string delimiter: the language uses it for loop
% labels, as Rust does, so treating it as a quote would swallow the rest of
% the line.
\lstdefinestyle{txhdl}{
  basicstyle=\ttfamily\listingsize,
  keywordstyle=\bfseries\color{kwblue},
  commentstyle=\itshape\color{cmtgreen},
  stringstyle=\color{strred},
  numberstyle=\tiny\color{numgray},
  morekeywords={mod,use,pub,const,type,struct,enum,trait,impl,fn,async,let,mut,
    match,if,else,loop,while,for,in,break,continue,return,as,await,
    module,bus,channel,transaction,protocol,pipeline,flow,seq,config,
    port,stage,pipe,instance,connect,bind,tag,domain,blackbox,foreign,
    property,role,signal,handler,true,false,
    sequence,interface,modport,implements,package,import,var,wire,func,proc,
    case,default,next,fork,branch,join,state,fsm},
  morecomment=[l]{//},
  morestring=[b]",
  showstringspaces=false,
  columns=fullflexible,
  keepspaces=true,
  breaklines=false,
  % Framed, so a listing is bounded rather than floating in the column.
  frame=single,
  rulecolor=\color{framegray},
  framesep=3pt,
  backgroundcolor=\color{bgshade},
  xleftmargin=3pt,
  xrightmargin=3pt,
  aboveskip=6pt,
  belowskip=6pt,
  % Every listing is numbered and captioned, so it can be referred to by
  % number. `caption` without `float` numbers the listing and leaves it
  % where it was written, which is what a listing in running prose needs.
  captionpos=b,
  abovecaptionskip=2pt,
  belowcaptionskip=6pt,
}

% Figures drawn as text. Framed the same way, and with the highlighting
% switched off, because a box-and-arrow diagram is not code and half its
% words turning blue reads as an accident. Setting a style adds keys rather
% than clearing the ones already set, so the three roles are neutralised by
% hand rather than by leaving them out.
\lstdefinestyle{ascii}{
  basicstyle=\ttfamily\listingsize,
  keywordstyle=\ttfamily,
  commentstyle=\ttfamily,
  stringstyle=\ttfamily,
  columns=fullflexible,
  keepspaces=true,
  frame=single,
  rulecolor=\color{framegray},
  framesep=3pt,
  backgroundcolor=\color{bgshade},
  xleftmargin=3pt,
  xrightmargin=3pt,
  aboveskip=6pt,
  belowskip=6pt,
}
\lstset{style=txhdl}

% House TikZ styles. `shorten` lives on the arrow styles rather than on each
% arrow, so every arrow in the document gets the gap, including ones added
% later. TikZ clips a path at a node's border, so without it an arrowhead
% butts against the box with no gap at all. Keep `node distance` well above
% twice the shorten, or the arrow shrinks to nothing.
\tikzset{
  box/.style={draw,rounded corners=2pt,align=center,inner sep=4pt,
              font=\footnotesize},
  boxf/.style={box,fill=bgshade},
  lbl/.style={font=\scriptsize\itshape,align=center},
  fl/.style={-{Stealth[length=4pt]},shorten >=2.5pt,shorten <=2.5pt},
  fld/.style={fl,dashed},
}


\usepackage{graphicx}

\InputIfFileExists{buildstamp.tex}{}{}
\providecommand{\buildstamp}{Draft build (outside the Bazel flow).}

\title{TxHDL: What Has Been Built}

\author{Filip Filmar and Dragiša Janković%
\thanks{Both authors are independent. Filip Filmar is the
corresponding author, at f@hdlfactory.com; Dragiša Janković is
reached at linkedin.com/in/dragisa-jankovic.}%
\thanks{Two pages on the whole of it, for a reader deciding whether
to read the rest. Every claim here is a target in the project
repository \code{HDL/txhdl} that the build runs; the cover document,
\code{//docs:cover}, names the fourteen documents that say more.}%
\thanks{The authors used a large language model, Claude, as an
assistant in exploring the concepts and in writing and constructing
the documents and the programs; every commit of the repository says
so and carries the prompts that led to it, verbatim.}%
\thanks{\buildstamp}}

\markboth{TxHDL: What Has Been Built}{Filmar and Janković: TxHDL: What Has Been Built}

\begin{document}
\maketitle

\section{The Rule}

TxHDL is a hardware description language that is a Rust library. It
was built by deleting every construct Rust already provides. A
transaction is a \code{struct}; a unit is a \code{struct} that
implements \code{Unit}; a pipeline is an \code{async fn}; a
configuration is a trait; a clock is a type. \code{rustc} is the
parser, and its type checker is the design rule checker: two drivers
on a wire are a move error, a clock domain crossed without a
\code{Crossing} is a type error, and a cycle counted in combinational
logic is impossible because a plain \code{fn} has no \code{.await} to
write in.

What is left is a unit, and a unit is a netlist. Listing~\ref{lst:w-unit}
is one whole: a counter that ticks once in eight cycles while
enabled. \code{\#[lower]} gives it a Verilog module and a VHDL entity
beside the Rust, and the build simulates both against the Rust run's
own trace, every time.

\lstinputlisting[rangebeginprefix=//\ begin\{,rangebeginsuffix=\},%
rangeendprefix=//\ end\{,rangeendsuffix=\},includerangemarker=false,%
linerange=unit-unit,caption={A unit, whole. Its netlist and its
waveform are on the cheat sheet, \code{//docs:cheatsheet}.},%
label={lst:w-unit}]{lib/examples/ex_cheat.rs}

\section{What Is In It}

\begin{table}[htbp]
\caption{The pieces, and where each is said in full.}
\label{tab:w-pieces}
\centering\footnotesize
\begin{tabular}{@{}lll@{}}
\toprule
Piece & What it is & Document \\
\midrule
The runtime & units, channels, clocks, time & \code{runtime} \\
The lowering & Verilog and VHDL from a unit & \code{embedding} \\
Parts & buffer, FIFO, station & \code{station} \\
The AXI link & five channels, two ends & \code{axi} \\
The AXI router & one host, N peripherals & \code{axi} \\
The network & hosts and peripherals on a lattice & \code{noc} \\
Vreteno & an RV32IM core, on a board & \code{vreteno} \\
The GPU & a rasteriser and a framebuffer & \code{gpu} \\
\bottomrule
\end{tabular}
\end{table}

\section{Vreteno, a Core}

Vreteno is an RV32IM core: fetch, execute and writeback, with the
multiply and divide as a sequence, the machine control registers, the
timer and the external interrupt. It is checked against a reference
model of the instruction set written as a program, in lockstep, every
cycle: the program counter, the thirty-one registers and the control
registers must agree after each retirement, on a demonstration
program and on sixty-four random ones.

It runs on hardware. The netlist goes through a hermetic Vivado onto
an Alinx AX7A200, and the board says on its serial port what the
simulation says it will: \code{OK}, and then whatever is typed at it,
echoed.

Its bus is AXI4, and every device behind it is an AXI peripheral,
with the router deciding by address which one a burst reaches.

\section{The AXI Link}

A channel carries one transaction. AXI is five channels that only
mean something together, so the library provides the link rather than
the channels: \code{axi()} makes every channel of it at once and
hands back two ends, the way \code{chan()} hands back a sender and a
receiver.

A host client issues a burst and awaits its answer; several may be
outstanding and they may be awaited in any order. A peripheral client
accepts a whole decoded transaction, with a write's beats already
gathered, and answers it from whatever process holds it. Neither
names a channel, counts a beat, watches \code{last} or writes an
identifier. Between the two ends are two units that do the protocol
and the five AXI channels, which is where the router goes; a router
changes no client code at all.

\section{A Network}

A router puts several peripherals behind one host. For several of
both, there is a lattice of nodes, each with a two-way link to its
four neighbours and a fifth, the exit, where a bridge with AXI on one
side and packets on the other attaches. Every link is two virtual
channels, one for requests and one for responses, which is what keeps
a full path of requests from holding up the answers that would empty
it. Routing is dimension order, which on a lattice needs nothing else
said to be free of deadlock.

Neither end knows the network is there. A host issues a burst and
awaits its answer as it would on a link; a peripheral accepts a
transaction and answers it as it would. What the network adds is the
renaming: an identifier belongs to the host that made it, so a
peripheral bridge gives every request one of its own and remembers
the node and the identifier to answer under. Two hosts on one memory
is the test that found that, and it failed before the renaming was
written.

\section{A GPU}

\begin{figure}[htbp]
\centering
\includegraphics[width=0.52\columnwidth]{scene.png}
\caption{Eight display list entries, drawn by the rasteriser, written
into memory over AXI, read back and written out as this file.}
\label{fig:w-scene}
\end{figure}

Figure~\ref{fig:w-scene} was drawn by hardware. A rasteriser walks
the box a display list entry gives it, one pixel a cycle, keeping
three edge functions incrementally, and writes each pixel that is
inside the primitive as a single-beat AXI burst; the framebuffer is a
memory behind the other end of the link. Both are lowered units, and
the build simulates both netlists against the run that drew a smaller
version of this picture.

The rasteriser was the first client of the AXI link that is hardware
rather than a testbench, and writing it is what showed that the
link's two client conveniences are simulation-side by construction
and that a unit holds the channel ends itself.

\section{Parts}

A part is hardware the language provides, written in its own lowerable
subset and checked as every unit is, so nothing in the library is a
primitive the compiler knows and the user cannot read. The channel
itself is one: an elastic buffer of two, with \code{valid} and
\code{ready} registered on both sides, which is what sits between two
lowered units wherever the simulation had a channel. A FIFO is
another, for a consumer that takes in bursts. A reservation station is
the third: it gathers the operands of an operation that arrive on
separate channels, out of order, and fires when the last of them has,
in two to ten inputs, each written out by a macro because the lowering
reads a body and not a loop over ports. The AXI link and its router
are parts too.

\begin{figure}[htbp]
\centering
\input{paper_cheat_timing.tex}
\caption{The unit of Listing~\ref{lst:w-unit}, running. The build drew
this from the trace that same run wrote, and generated a testbench
from it that replays the netlist.}
\label{fig:w-wave}
\end{figure}

\section{From an Empty Directory}

A reader with nothing installed but Bazelisk writes four files, and
the fourth is a blinky that builds, runs, and lowers to Verilog and
VHDL. That workspace is kept in the tree and built as a check, so the
tutorial cannot rot: it depends on the project the way anyone else
would, through the public mirror, and the build fetches the toolchain,
the simulators, the waveform tools and a TeX installation itself.
Nothing is installed by hand, on a workstation or in CI.

\section{How It Is Checked}

Nothing here is checked by reading it. Every example is a program the
build runs, and its output is included in a document from the run, so
a document cannot drift from the code. Every lowered unit is
simulated twice more, as VHDL under \code{nvc} and as Verilog under
Verilator, against the trace the Rust run wrote, at its ports; that
is what the sixty-two tests under \code{//docs} are. Every claim
about what Rust accepts was compiled first, as a probe kept in the
tree, including the ones whose answer was no.

\section{What Is Not Done}

The list is short and it is kept in the open. The lowering emits a
name as it is written and neither escapes nor rejects the reserved
words of its targets, so a register called \code{next} produces VHDL
that will not analyse; three names in the GPU had to change for that
reason. A unit's ports are a positional tuple and not a named struct,
which a seven-port unit feels. The board's top level is still wired by
hand and is out of date since the core moved to AXI. The GPU has no
depth buffer, no interpolation and no texture. Fuzzing the AXI link is
not yet a step of the build, and the network has its nodes, its
bridges and its tests but not yet the four-cornered demonstration it
was asked for. Each of those is an issue on the
repository rather than a paragraph here.

\begin{table}[htbp]
\caption{The tree, in lines of Rust and in targets.}
\label{tab:w-numbers}
\centering\footnotesize
\begin{tabular}{@{}lr@{}}
\toprule
The runtime and its macros & 8\,544 \\
Parts: the link, the router, the network & 3\,915 \\
Examples, forty-two of them & 4\,335 \\
Vreteno & 3\,738 \\
The GPU & 1\,230 \\
\midrule
Tests, of which 62 are co-simulations & 71 \\
Documents & 16 \\
\bottomrule
\end{tabular}
\end{table}

\end{document}
